\documentclass[11pt]{article}
\usepackage[margin=1in]{geometry}
\usepackage{amsmath,amssymb,amsthm,mathtools}
\usepackage[hidelinks]{hyperref}
\newtheorem{theorem}{Theorem}
\newtheorem{corollary}[theorem]{Corollary}
\newtheorem{lemma}[theorem]{Lemma}
\title{Cyclic Wronskians, factorial discriminants\\and positive auxiliary factorizations}
\author{Internal proof and exact replay; no priority claim}
\date{17 November 2026}
\begin{document}
\maketitle
For every integer gap $d\ge1$, define the monic auxiliary by
\[
 B_{0,\gamma,d}=1,\qquad
 B_{n+1,\gamma,d}=(z+dn+\gamma+d-1)B_{n,\gamma,d}
                                      -z\partial_zB_{n,\gamma,d}.
                                                               \tag{1}
\]
At the physical parameter $\gamma=n+b-1$ put
$\mathcal B_n(z;b,d)=B_{n,n+b-1,d}(z)$ and let
\[
 R_{n,k}(b,d)=[z^k]\mathcal B_n(z;b,d),\qquad n,k\ge0.
\]
Thus $R$ is lower triangular with diagonal entries one.
All generating identities below are formal power-series identities
in characteristic zero and hold for every real $b$; positivity is not
assumed in deriving them. The counterexamples use the physical
ordinary-factorial shape $b=1$.

Theorem~\ref{thm:discriminant} below identifies a generalized Hurwitz
determinant with a normalized factorial-moment discriminant at every gap.
It proves nonnegativity of that determinant for every positive Gamma shape,
and of the last three special minors at gap three. Strict nonvanishing
and total nonnegativity of all generalized Hurwitz minors remain open.

\section{An exact all-gap Riordan representation}
\begin{theorem}\label{thm:Riordan}
Set
\[
 P_d(u)=\frac{(1-u)-(1-u)^{d+1}}d,\qquad P_d(0)=0,
 \quad P_d'(0)=1,
\]
and let $f_d(t)$ be its formal compositional inverse with $f_d(0)=0$.
Then
\[
 \sum_{n\ge0}\mathcal B_n(z;b,d)\frac{t^n}{n!}
 =g_{b,d}(t)e^{zf_d(t)},\qquad
 g_{b,d}(t)=f_d'(t)(1-f_d(t))^{2-b}.                         \tag{2}
\]
Equivalently, $w=1-f_d(t)$ satisfies
$t=(w-w^{d+1})/d$ and
\[
 g_{b,d}(t)=\frac{d\,w^{2-b}}{(d+1)w^d-1}.
\]
The inverse matrix has the explicit coefficient formula
\[
 (R^{-1})_{n,k}=\frac{n!}{k!}[u^n]
          P_d'(u)(1-u)^{b-2}P_d(u)^k.                      \tag{3}
\]
\end{theorem}
\begin{proof}
The fixed-$\gamma$ generating function obtained directly from (1) is
\[
 \sum_{n\ge0}B_{n,\gamma,d}(z)\frac{s^n}{n!}
 =(1-ds)^{-(\gamma+d-1)/d}
            e^{z(1-(1-ds)^{1/d})}.
\]
Indeed its derivative satisfies
$(1-ds)F_s=(z+\gamma+d-1)F-zF_z$, and its constant term is one.
Coefficient extraction at $\gamma=n+b-1$ gives
\[
 \frac{\mathcal B_n}{n!}
 =d^n\operatorname{res}_{v=0}
  (1-v)^{-(n+b+d-2)/d}
  e^{z(1-(1-v)^{1/d})}\frac{dv}{v^{n+1}}.
\]
Make the formal invertible substitution $u=1-(1-v)^{1/d}$.
Since $dv=d(1-u)^{d-1}du$ and
$v(1-u)=dP_d(u)$, the last expression becomes
\[
 \operatorname{res}_{u=0}
       \frac{(1-u)^{2-b}e^{zu}}{P_d(u)^{n+1}}\,du.
\]
Substituting $t=P_d(u)$ proves (2) by the formal change-of-variable
formula for residues. No analytic contour or convergence assertion
is used.

For completeness the elementary multiplication law for coefficient
arrays with column generating functions $g(t)f(t)^k/k!$ is
\[
 (g,f)(h,q)=(g(t)h(f(t)),q(f(t))).
\]
It follows by substituting the column generating functions and
comparing coefficients. Taking $q=P_d$ and
$h(u)=1/g(P_d(u))=P_d'(u)(1-u)^{b-2}$ gives the identity matrix,
proving (3).
\end{proof}

The identity is compatible with the auxiliary construction in the
\href{HIGHER_GAP_ADJOINT_RANK.tex}{higher-gap adjoint note};
its proof here uses only (1).
For gap three,
$P_3(u)=u-2u^2+\tfrac43u^3-\tfrac13u^4$.
This inverse-polynomial structure gives a possible entry point for
factorizations, but the entire auxiliary array is not totally
nonnegative.

\section{A positive factorization with the shape held fixed}
The row-dependent shape in $R$ is essential. With $\gamma$ held fixed,
write $\mathcal R_{m,k}(\gamma,d)=[z^k]B_{m,\gamma,d}(z)$.
For zero-based indices $0\le i,j\le N$, define
\[
 D_N=\operatorname{diag}(0!,1!,\ldots,N!),\quad
 (T_{d,N})_{i,j}=\begin{cases}\binom{i-j+d-1}{d-1},&i\ge j,\\0,&i<j,
 \end{cases}
\]
and the upper bidiagonal matrix
\[
 (U_{\gamma,d,N})_{j,j}=\gamma+d-1-j,\qquad
 (U_{\gamma,d,N})_{j-1,j}=j.
\]
\begin{theorem}\label{thm:fixed}
The leading production block of the fixed-shape array factors exactly as
\[
 \Pi^{(\gamma,d)}_N=D_N T_{d,N}U_{\gamma,d,N}D_N^{-1}.          \tag{4}
\]
If $\gamma+d-1\ge N$, then both this production block and the
leading block $\mathcal R[0..N,0..N]$ are totally nonnegative.
In particular, at gap three with $\gamma=n+b-1$ and $b\ge0$,
the fixed-shape auxiliary array through row and column $n+1$ is
entirely totally nonnegative. More generally, for $d\ge2$ the block
through $n+d-2$ has this property throughout the physical range.
\end{theorem}
\begin{proof}
The fixed-shape generating function in the proof of
Theorem~\ref{thm:Riordan} gives production data
\[
 A(u)=(1-u)^{1-d},\qquad Z(u)=(\gamma+d-1)(1-u)^{-d}.
\]
If $h_j=[u^j](1-u)^{-d}$ and $c=\gamma+d-1$, the factorial-scaled
production entry is
\[
 c h_{i-j}+j(h_{i-j+1}-h_{i-j})
      =(c-j)h_{i-j}+j h_{i-j+1}.
\]
This is exactly the $(i,j)$ entry of $T_{d,N}U_{\gamma,d,N}$,
proving (4), including its superdiagonal boundary entries.

Let $G_N$ be the lower triangular matrix of ones. Then
$T_{d,N}=G_N^d$, by the coefficient identity for $(1-u)^{-d}$.
Also
\[
 G_N=(I+E_{N,N-1})\cdots(I+E_{1,0}).
\]
Every factor is totally nonnegative. If $c\ge N$, the matrix $U$
is a nonnegative upper bidiagonal matrix, hence is totally
nonnegative as well. Positive diagonal scaling and the finite
Cauchy--Binet identity establish the assertion about $\Pi_N$.

For the coefficient array use induction on $N$.
Let $Q_N$ have first row $(1,0,\ldots,0)$ and remaining rows
$\Pi^{(\gamma,d)}[0..N-1,0..N]$. Every minor of $Q_N$ either
is a production minor, is zero, or expands along its first row
to a production minor with positive sign. Thus $Q_N$ is totally
nonnegative. The production identity gives
\[
 \mathcal R_N=\operatorname{diag}(1,\mathcal R_{N-1})Q_N.
\]
Induction and Cauchy--Binet prove total nonnegativity of
$\mathcal R_N$. For the stated physical block
$N=n+d-2$, its parameter satisfies $c=N+b\ge N$.
\end{proof}

This is a finite positive factorization at an unbounded set of
indices and parameters, rather than evidence from selected minors.
It supplies an actual planar-network construction for the fixed-shape
auxiliary blocks. It does not yet supply a network for the cyclic
Wronskian determinants.

\section{The exact cyclic Wronskian transfer}
Now fix $d=3$, $B=B_{n,\gamma,3}$ and $C=B_{n+1,\gamma,3}$.
Write their cyclic sections as
\[
 B(z)=A_0(z^3)+zA_1(z^3)+z^2A_2(z^3),\qquad
 C(z)=C_0(z^3)+zC_1(z^3)+z^2C_2(z^3).
\]
The adjacent-row recurrence (1) gives the actual Wronskian of the
\href{gap_three_hurwitz.tex}{generalized Hurwitz target} as
\[
 W(z)=\frac{C(\alpha z)B(\bar\alpha z)-B(\alpha z)C(\bar\alpha z)}
              {(\alpha-\bar\alpha)z},\qquad \alpha=e^{i\pi/3}.
\]
The quotient is polynomial at zero.
\begin{theorem}\label{thm:transfer}
If $W(z)=W_0(z^3)+zW_1(z^3)+z^2W_2(z^3)$, then
\begin{align*}
 W_0(v)&=(A_0C_1-A_1C_0)(-v),\\
 W_1(v)&=(A_0C_2-A_2C_0)(-v),\\
 W_2(v)&=(A_1C_2-A_2C_1)(-v).                              \tag{5}
\end{align*}
Consequently the common-section-root obstruction is precisely a
rank drop of the two cyclic-section rows
\[
 \begin{pmatrix}A_0(u)&A_1(u)&A_2(u)\\C_0(u)&C_1(u)&C_2(u)\end{pmatrix}.
\]
If $B=\sum_j b_jz^j$ and $C=\sum_j c_jz^j$, the same transfer is
\[
 W(z)=\sum_{j<k}\chi_{k-j}(b_jc_k-b_kc_j)z^{j+k-1},\qquad
 \chi_\ell=\frac{\sin(\ell\pi/3)}{\sin(\pi/3)}.              \tag{6}
\]
The weights $\chi_1,\chi_2,\ldots$ are periodic:
$1,1,0,-1,-1,0,\ldots$.
\end{theorem}
\begin{proof}
The raising identity is $C=(z+3n+\gamma+2)B-zB'$.
Substituting it into the divided difference gives exactly the defining
rotated Wronskian, since $\alpha-\bar\alpha=i\sqrt3$.
Next use $\alpha^3=\bar\alpha^3=-1$ and
$\alpha\bar\alpha=\alpha+\bar\alpha=1$.
Insert the three cyclic sections of $B,C$ into the numerator and
combine its three pairs of distinct residue classes.
The pairs $(0,1),(0,2),(1,2)$ contribute respectively the three
minors in (5), multiplied by $1,z,z^2$ and evaluated at $u=-z^3$.
This proves (5); their simultaneous vanishing is exactly rank at
most one of the displayed two-row matrix.
Alternatively, collect coefficient pairs $j<k$ directly in the
rotated numerator. Their multiplier is
\[
 \frac{\alpha^k\bar\alpha^j-\alpha^j\bar\alpha^k}{\alpha-\bar\alpha}
       =\frac{\sin((k-j)\pi/3)}{\sin(\pi/3)},
\]
proving (6).
\end{proof}

At the physical parameter, Theorem~\ref{thm:fixed} makes the
adjacent auxiliary coefficient minors $b_jc_k-b_kc_j$ nonnegative.
Formula (6) shows explicitly why that does not make the transferred
Wronskian minors positive term by term: the cyclic weights have both
signs. The uniform $H_3(W)$ task therefore asks for control of the
three specific Pl\"ucker minors in (5), using the order-three
Toeplitz factor of (4). The weaker simplicity route needs only that
these three minors have no common nonzero root, by the existing
\href{HIGHER_GAP_ADJOINT_RANK.tex}{adjoint rank criterion}.
No such uniform transfer is proved here.

\subsection*{Two uniform tail determinant identities}
There is an additional consequence of the cyclic Pl\"ucker structure
that holds without a positivity assumption. Write
\[
 F(t)=t^{2n}W(1/t)=\sum_{j=0}^{2n}w_jt^j,\qquad w_0=1,
\]
with $w_j=0$ outside $0\le j\le2n$. Define the special minors by
\[
 \Delta_p(W)=\det(w_{3j-i})_{{k\le i\le k+r-1}\atop{1\le j\le r}},
 \quad r=\lceil p/2\rceil,\quad k=2r-p+1,
 \qquad \Delta_0(W)=1.
\]
These are the same descending-coefficient generalized Hurwitz minors
as in the uniform target.
\begin{theorem}\label{thm:tail}
For every pair of monic polynomials $B,C$ of consecutive degrees
$n,n+1$, $n\ge1$, the divided difference of Theorem~\ref{thm:transfer}
satisfies
\[
 \Delta_{2n-1}(W)=B(0)\Delta_{2n-2}(W),\qquad
 \Delta_{2n}(W)=W(0)\Delta_{2n-2}(W).                       \tag{7}
\]
The identities remain valid when $B(0)$, $W(0)$ or
$\Delta_{2n-2}(W)$ is zero.
\end{theorem}
\begin{proof}
Put $D_{ij}=A_iC_j-A_jC_i$.
The Pl\"ucker identity
$A_0D_{12}-A_1D_{02}+A_2D_{01}=0$ shows that $B(-z)W(z)$
has no powers congruent to two modulo three: its residue-two section
is exactly that expression, evaluated at $-z^3$.
Reverse this degree-$3n$ product and put
\[
 \widehat B(t)=(-1)^nt^nB(-1/t)
       =\sum_{j=0}^n\widehat b_jt^j.
\]
Then $\widehat b_0=1$, $\widehat b_n=(-1)^nB(0)$, and
$\widehat B(t)F(t)$ has no powers congruent to one modulo three.
For $j=1,\ldots,n$, its coefficient at $t^{3j-2}$ is zero.
If
\[
 K_i=(w_{3j-i})_{j=1}^n,\qquad 2\le i\le n+2,
\]
these identities give the exact row relation
\[
 K_2+\sum_{i=3}^{n+2}\widehat b_{i-2}K_i=0.
\]
Expand $\Delta_{2n-1}(W)=\det(K_2,K_3,\ldots,K_{n+1})$
using this relation. All terms except the last have repeated rows,
so
\[
 \Delta_{2n-1}(W)=(-1)^n\widehat b_n
                   \det(K_3,\ldots,K_{n+2}).
\]
The last determinant has first column $(w_0,0,\ldots,0)^T$.
Expanding along it and reindexing the remaining rows and columns
gives $\Delta_{2n-2}(W)$. This proves the first identity.
For $n=1$ the remaining determinant is empty and equals one,
so the same proof includes that case.

For the even minor, its last column has only the bottom nonzero
entry $w_{2n}=W(0)$, by the degree bound on $F$.
Expansion along that column proves the second identity.
No step divides by a constant coefficient, determinant or rank
assumption.
\end{proof}

In the physical family, $B(0)>0$ and $W(0)>0$.
Consequently the last two special minors have the sign of
$\Delta_{2n-2}(W)$. A proof using strict positivity of all special
minors therefore has independent positivity obligations only through
index $2n-2$. Their uniform positivity or nonvanishing remains open.

\subsection*{Why the auxiliary-block property alone does not transfer}
Take $C_0(z)=1$ and
\[
 C_{m+1}(z)=(z+9)C_m(z)-zC_m'(z),\qquad 0\le m<9.
\]
The coefficient array through degree nine has upper bidiagonal
production matrix with diagonal entries $9-j$ and superdiagonal
entries one. Its coefficients therefore form a totally nonnegative
block by the same induction. Form the gap-three Wronskian from $C_8$:
\[
 V(z)=C_8(\alpha z)C_8(\bar\alpha z)
 -\frac{\alpha C_8'(\alpha z)C_8(\bar\alpha z)
        -\bar\alpha C_8'(\bar\alpha z)C_8(\alpha z)}{i\sqrt3},
 \qquad \alpha=e^{i\pi/3}.
\]
Writing $V=\sum_{j=0}^{16}v_jz^{16-j}$, its generalized Hurwitz
minor in rows $1,2,3$ and columns $1,2,3$ is
\[
 \det\begin{pmatrix}
 870&1481382&1821867821\\
 44&127869&181136695\\
 1&11256&16846286
 \end{pmatrix}=-60115263786123<0.
\]
This auxiliary is outside the required arithmetic progression:
the raising constants are all nine, instead of increasing by three.
The example excludes a transfer theorem based only on total
nonnegativity of the auxiliary coefficient block, even with the
same first-order raising-operator form. A valid transfer must use
the order-three Toeplitz factor in (4), or another part of the
actual parameter compatibility.

\section{A uniform discriminant identity at every gap}
For integers $n,d\ge1$, put
\[
 T_{n,\gamma,d}(t)=\sum_{k=0}^n\binom nk(\gamma)_{dk}t^k,
 \qquad (\gamma)_j=\prod_{r=0}^{j-1}(\gamma+r),\quad(\gamma)_0=1.
\]
Throughout this section, $\operatorname{Disc}_tT_{n,\gamma,d}$ means the
\emph{fixed-degree-$n$ coefficient discriminant}: compute it with
indeterminate $\gamma$ and then specialize that polynomial.
At a negative integer the actual specialized degree of $T$ may drop;
recomputing the discriminant at that smaller degree is a different
operation. There is no ambiguity for $\gamma>0$.

For $0\le s\le d-2$, let $P_s=B_{n+s,\gamma,d}$ and write
\[
 P_s(z)=\sum_{j=0}^{d-1}z^jM_{s,j}(z^d),\qquad
 D_j(u)=\det M_{\widehat j}(u).
\]
Here $M_{\widehat j}$ deletes column $j$ from the $(d-1)$-by-$d$
cyclic-section matrix. Define
\[
 W^{[d]}_{n,\gamma}(z)=
    \sum_{j=0}^{d-1}z^{d-1-j}D_j((-1)^dz^d).
\]
At $d=1$ the row set is empty, its determinant is one, and $W^{[1]}=1$.
The polynomial $W^{[d]}$ is monic of degree $(d-1)n$; this is proved
below. Write its descending coefficients as
\[
 z^{(d-1)n}W^{[d]}_{n,\gamma}(1/z)=\sum_jw_jz^j,
 \qquad
 \mathcal D_{n,d}(\gamma)=\det(w_{dj-i})_{i,j=1}^{n-1},
\]
with $w_j=0$ outside the polynomial's range and
$\mathcal D_{1,d}=1$.
For $d=3$ this $W^{[3]}$ is exactly the earlier $W$, and
$\mathcal D_{n,3}=\Delta_{2n-2}(W)$.

\begin{theorem}\label{thm:discriminant}
For every pair of integers $n,d\ge1$, the following is a polynomial
identity in $\gamma\in\mathbb C$, with $N=\binom n2$:
\[
 \operatorname{Disc}_t T_{n,\gamma,d}
 =(-1)^N d^N
   \left(\prod_{k=1}^n k^k\right)
   \left(\prod_{k=1}^{n-1}(\gamma)_{dk}\right)
        \mathcal D_{n,d}(\gamma).                           \tag{D}
\]
In particular, for $\gamma>0$, $T_{n,\gamma,d}$ is squarefree if and
only if $\mathcal D_{n,d}(\gamma)\ne0$.
\end{theorem}
The proof identifies the relevant schemes, including multiplicities.
A comparison of zero sets alone would not establish this identity.

\subsection*{The bounded cyclic presentation}
\begin{lemma}\label{lem:general-fitting}
Let $A$ be any commutative ring with identity and let
$P_0,\ldots,P_{d-2}\in A[z]$ be monic of consecutive degrees
$n,\ldots,n+d-2$. Form $M,D_j,W,w_j$ exactly as above, and put
$E_j=(-1)^jD_j$.
Then $W$ is monic of degree $(d-1)n$, and, up to signed changes
of free bases, there is an exact sequence of $A$-modules
\[
 A^{n-1}\xrightarrow{H^T}A^{n-1}\longrightarrow
 A[u]/(E_0,\ldots,E_{d-1})\longrightarrow0,
 \qquad H=(w_{dj-i})_{i,j=1}^{n-1}.
\]
Consequently the quotient's zeroth Fitting ideal is $(\det H)$.
Over a discrete valuation ring, if $\det H\ne0$, its length equals
$v(\det H)$. No genericity or nonzero-constant assumption is required.
\end{lemma}
\begin{proof}
For $d=1$ the quotient is zero and $H$ is an identity matrix, so the
assertion is direct. Suppose $d\ge2$.
Assign weight $de+j$ to $u^ee_j$ in a multiplier row.
The cyclic row $M_s$ has unique greatest-weight monomial
$u^{q_s}e_{r_s}$, with coefficient one, where
\[
 q_s=\lfloor(n+s)/d\rfloor,\qquad r_s=(n+s)\bmod d.
\]
The $d-1$ pivots are distinct. Write $n-1=da+l$, $0\le l<d$;
the missing column is $l$, and $Q=\sum_sq_s=(d-1)a+l$.
The missing cofactor $E_l$ has degree $Q$ and leading coefficient
$1$ or $-1$: its unique maximum-weight determinant term is the
product of the distinct row pivots. More generally,
\[
 \deg E_j\le\left\lfloor\frac{(d-1)(n-1)+j}{d}\right\rfloor
 =\begin{cases}Q-1,&j<l,\\Q,&j\ge l.\end{cases}
\]
The unsigned minor omitting $l$ has leading pivot-permutation sign
$(-1)^{l(d-1-l)}$. In $W$, after $u=(-1)^dz^d$, its leading sign
is $(-1)^{dQ+l(d-1-l)}=1$; the exponent is congruent to $l(l+1)$
modulo two. Every other summand has lower degree. Thus $W$ is monic
of degree $(d-1)n$.

Set $t=n-2$ and $V=A[u]_{\le t}$, interpreted as zero when $t=-1$.
Allow multiplier bounds
\[
 b_j=\left\lfloor\frac{t-j}{d}\right\rfloor
 =\begin{cases}a,&j<l,\\a-1,&j\ge l.\end{cases}
\]
Negative bounds mean zero modules. The bounded multiplication map
\[
 \Phi:\bigoplus_{j=0}^{d-1}A[u]_{\le b_j}\longrightarrow V,
 \qquad(f_j)\longmapsto\sum_jf_jE_j
\]
has equal source and target ranks $n-1$.
Since $E_l$ is monic up to sign of degree $Q\le n-1$, every quotient
class has a representative in $V$. At $n=1$, $E_l$ is a unit and
the quotient is zero directly.

To prove exactness, let $P=\sum f_jE_j$ have degree at most $t$.
The row syzygies are $\sum_jM_{s,j}E_j=0$.
If a pivot component $f_{r_s}$ has a term $cu^e$ with $e\ge q_s$,
subtract $cu^{e-q_s}M_s$ from the multiplier vector.
This cancels the chosen term and introduces only monomials of
strictly smaller weight. Always choose an eligible term of maximal
weight. The process terminates, with
$\deg f_{r_s}<q_s$, exactly the required bounds on all pivot columns.
Their products with $E_j$ have degree at most $t$.
Monicity of $E_l$ then forces
$\deg f_l\le t-Q=a-1=b_l$ as well.
All reductions use leading coefficients $1$ or $-1$; the argument
also holds over rings with zero divisors.
Thus $\ker(V\to A[u]/(E))=\operatorname{im}\Phi$.

It remains to identify its matrix. Put $L=(d-1)n$, $m=n-1$ and
write $W=\sum_{k=0}^{d-1}z^kW_k(z^d)$. The cofactor definition gives
$W_k(v)=(-1)^{d-1-k}E_{d-1-k}((-1)^dv)$.
Use descending coefficient order in the target $V$.
For $1\le i\le m$, set
\[
 k_i=(L+i)\bmod d,\quad c_i=\lfloor(L+i)/d\rfloor,
 \quad\ell_i=m-c_i,\quad j_i=d-1-k_i.
\]
Then
\[
 (H^T)_{ji}=w_{dj-i}
 =[v^{m-j}]\,v^{\ell_i}W_{k_i}(v),\qquad
 d\ell_i+j_i=n-i-1.
\]
These columns enumerate exactly the bounded multipliers, once each.
The substitution $u=(-1)^dv$ and the displayed cofactor signs change
only signs in the free bases. Hence the matrix is $H^T$ up to signed
permutations. The Fitting assertion follows from this square
presentation; over a discrete valuation ring, Smith normal form
gives the stated length.
\end{proof}

\subsection*{The adjoint Wronskians and critical ideals}
Let $Q=1-D_z$ and let $\omega_0,\ldots,\omega_{d-1}$ be the $d$th
roots of unity. The polynomial just defined also has the expression
\[
 W^{[d]}_{n,\gamma}(z)=\frac1{\kappa_d}
  \det\bigl[\omega_j^k(Q^kB_{n,\gamma,d})(-\omega_jz)\bigr]
       _{1\le j\le d-1,\ 0\le k\le d-2},
\]
where $\omega_0=1$ and
$\kappa_d=\prod_{1\le i<j\le d-1}(\omega_j-\omega_i)$;
empty determinants and products at $d=1$ equal one.
To check the normalization, put $X=zQ$ and $c=\gamma+d-1$.
The recurrence gives
$B_{n+s}=\prod_{r=0}^{s-1}(X+dn+c+dr)B_n$.
Also $X^s=z^sQ^s$ plus a constant linear combination of
$z^kQ^k$, $k<s$.
Thus the adjacent-row evaluation determinant differs from the
adjoint minor by the factor $(-z)^{(d-1)(d-2)/2}$.
Expanding that evaluation determinant through the cyclic sections,
the Fourier minor omitting column $j$ has ratio
$e_{d-1-j}(\omega_1,\ldots,\omega_{d-1})=(-1)^{d-1-j}$ to
$\kappa_d$, because the nontrivial roots have polynomial
$1+x+\cdots+x^{d-1}$. These signs cancel the powers of $-z$,
giving exactly the unsigned cofactor formula above, with no extra
power of $z$ or $u$.

Put
\[
 R(z,\gamma)=z^{dn}T_{n,\gamma,d}((-1)^{d+1}z^{-d}),\qquad
 M=\gamma+dn.
\]
This is monic of degree $dn$ and has constant coefficient
$(-1)^{n(d+1)}(\gamma)_{dn}$.
The adjoint construction of the
\href{HIGHER_GAP_ADJOINT_RANK.tex}{higher-gap note} gives $d$ solutions
\[
 f_\omega(z)=e^{-\omega z}z^{1-M}B_{n,\gamma,d}(\omega z)
\]
of an order-$d+1$ adjoint equation ordinary at every $z\ne0$.
One may take the formal adjoint of
\[
 L=z^{-d-1}\left[\delta\prod_{\ell=0}^{d-1}(\delta-M-\ell)
                    -z^d(\delta-dn)\right],\qquad\delta=zD_z;
\]
its leading coefficient is a nonzero constant and is normalized to one.
The same construction gives
\[
 E=\det(f_\omega^{(k)})_{0\le k\le d-1}
       =\text{a local unit}\cdot R.
\]
These algebraic identities hold for every complex $\gamma$.
Indeed the gauge equation and the stripped determinant are polynomial
identities in $\gamma$, so their positive-shape derivations extend
coefficientwise to all complex values. The same formulas also hold
at $d=1$ by the identical coefficient calculation.
The solutions remain independent: after removal of the common power
of $z$, they are exponential polynomials with distinct exponents and
monic degree-$n$ polynomial factors.
Consequently their $d$-by-$(d+1)$ jet matrix has rank $d$ at every
nonzero point, by uniqueness of ordinary-point initial data.
Local branches of $z^{1-M}$ only multiply rows and triangular jet
gauges by units.

\begin{lemma}\label{lem:critical-ideal}
At every complex parameter and every $z\ne0$, the ideal generated
by $R,R'$ equals, up to invertible changes of generators, the ideal
of maximal minors of the first $d-1$ adjoint-jet columns.
After removing gauges and applying the cyclic Fourier transform,
this is the ideal $(D_0,\ldots,D_{d-1})$ at $u=z^d$.
\end{lemma}
\begin{proof}
Write the full jet columns as $C_0,\ldots,C_d$.
Wronskian differentiation gives
$E'=\det(C_0,\ldots,C_{d-2},C_d)$.
The ideals $(E,E')$ and $(R,R')$ agree because $E$ is a unit times $R$.
At a common zero of $E,E'$, the first $d-1$ columns must have rank
at most $d-2$; otherwise both final columns lie in their span and
the full jet rank drops. The full rank also makes their rank at
least $d-2$. Choose those independent $d-2$ columns together with
$C_{d-1},C_d$ as a local basis. The remaining early column has two
coordinates outside their early span. The two determinants $E,E'$
are these two coordinates up to units, and the maximal early minors
generate the same coordinate ideal. This proves equality of ideals,
not merely equality of their zero sets. Outside their common zero
set both local ideals are units. For $d=1$ both ideals are
units, since a solution of an ordinary second-order equation cannot
have both its value and derivative zero.

Removing the gauges leaves columns
$\omega^k(Q^kB)(\omega z)$, $0\le k\le d-2$.
The polynomials $B_n,\ldots,B_{n+d-2}$ are unit triangular
combinations of $(zQ)^kB_n$, and their column transformation has
powers of $z$ on the diagonal, invertible when $z\ne0$.
The discrete Fourier matrix, followed by removal of the factors
$1,z,\ldots,z^{d-1}$, gives the cyclic rows $M$.
Every operation is invertible locally, proving the last assertion.
\end{proof}

\begin{lemma}\label{lem:no-origin}
For every complex $\gamma$, the cyclic matrix $M(0)$ has rank $d-1$.
\end{lemma}
\begin{proof}
The case $d=1$ is the empty-minor convention.
For $d\ge2$ put $g=e^{-z}B_{n,\gamma,d}(z)$.
Its first $d$ Taylor derivatives are
\[
 g^{(k)}(0)=(-1)^k(Q^kB)(0)
   =(-1)^k\prod_{j=0}^{n-1}(\gamma-k+d-1+dj),\quad0\le k<d.
\]
The product polynomials have pairwise disjoint root sets, so at most
one of these derivatives is zero.
For the $d-1$ functions $g(-\omega_jz)$ with nontrivial $\omega_j$,
every Fourier minor using $d-1$ of the first $d$ Taylor indices is
nonzero after omission of a possible zero index; the elementary
symmetric factors of the nontrivial roots are all $1$ or $-1$.
The first available $d-1$ Taylor indices give the unique lowest-order
Wronskian term, whose coefficient is nonzero.
Its order is the sum of those indices minus $(d-1)(d-2)/2$,
hence is at most $d-1$.
This Wronskian is a unit times $W^{[d]}$, by the root-of-unity formula.
If all its $d$ cyclic-section constants vanished, its order would be
at least $d$, a contradiction. Thus not all $D_j(0)$ vanish.
\end{proof}

\subsection*{The two leading coefficients}
Write $D_{n,d}(\gamma)=\operatorname{Disc}_tT_{n,\gamma,d}$.
For fixed $n,d$, with $N=\binom n2$, we claim
\[
 \mathcal D_{n,d}(\gamma)=d^N\gamma^{(d-1)N}
            +\text{lower powers},
\qquad
 D_{n,d}(\gamma)=(-1)^Nd^{2N}
       \left(\prod_{k=1}^n k^k\right)\gamma^{(2d-1)N}
            +\text{lower powers}.                          \tag{L}
\]
In particular neither polynomial is identically zero.

First, $\gamma^{-n}B_{n,\gamma,d}(\gamma z)\to(1+z)^n$
coefficientwise. The adjoint-minor formula therefore gives
\[
 \gamma^{-(d-1)n}W^{[d]}_{n,\gamma}(\gamma z)
       \longrightarrow(1+z+\cdots+z^{d-1})^n.
\]
The weighted degree of the determinant is $(d-1)N$.
Its leading coefficient is the special determinant of this limit.
Multiplying any coefficient series by a unit series in $t^d$ leaves
that determinant unchanged, since it is a unit upper triangular
column convolution. Thus replace
\[
 (1+t+\cdots+t^{d-1})^n=(1-t^d)^n(1-t)^{-n}
 \quad\text{by}\quad(1-t)^{-n}.
\]
Its coefficients are $h_j=\binom{n+j-1}{n-1}$, zero for negative
indices. Reverse and transpose the determinant and append the last
unit row and column. The resulting $n$-square determinant is
\[
 \det\left[\binom{x_i+j-1}{n-1}\right]_{i,j=1}^n,
 \qquad x_i=d(n-i).
\]
Successive column finite differences give polynomial degrees
$n-1,n-2,\ldots,0$, so this is a constant times the Vandermonde
in the $x_i$. At $x_i=n-i$ its value is one, by the corresponding
triangular coefficient matrix. Scaling all $x_i$ by $d$ multiplies
it by $d^N$. This proves the first formula in (L).

For the second, let $Y_\gamma$ have unit-rate Gamma shape $\gamma$
and put $X_\gamma=Y_\gamma^d$. The monic reversed moment polynomial
is $P_n(x)=\mathbb E(x+X_\gamma)^n$; reversal preserves its
discriminant, so $\operatorname{Disc}P_n=D_{n,d}$.
Center by $\mathbb EX_\gamma$ and scale by
$s_\gamma=d\gamma^{d-1/2}$.
The Gamma cumulants
$\kappa_r(Y_\gamma)=\gamma(r-1)!$ show that every fixed moment of
$(Y_\gamma-\gamma)/\sqrt\gamma$ tends to its standard Gaussian
moment. Expanding the integer power $Y_\gamma^d$ then gives the same
moment convergence for
$(X_\gamma-\mathbb EX_\gamma)/s_\gamma$.
Consequently
\[
 s_\gamma^{-n}P_n(s_\gamma x-\mathbb EX_\gamma)
 \longrightarrow H_n(x)=\mathbb E(x+G)^n,
\]
where $G$ is standard Gaussian.
The recurrence $H_n=xH_{n-1}+(n-1)H_{n-2}$ and
$H_n'=nH_{n-1}$ give an explicit resultant induction: if
$J_n=\operatorname{Res}(H_n,H_{n-1})$, then
$J_n=(n-1)^{n-1}J_{n-1}$ and $J_1=1$.
Thus $\operatorname{Disc}H_n=(-1)^Nn^nJ_n
=(-1)^N\prod_{k=1}^nk^k$.
Scaling discriminants proves the second leading coefficient in (L).
These are polynomial leading-term calculations; no assertion about
zeros follows merely from a numerical normal approximation.

\subsection*{Valuations, including every degree-drop parameter}
Fix $\gamma_0\in\mathbb C$ and work over the complete discrete
valuation ring $A=\mathbb C[[\gamma-\gamma_0]]$.
Lemma~\ref{lem:general-fitting} identifies the length of the cyclic
quotient with $v(\mathcal D_{n,d})$.
By Lemma~\ref{lem:no-origin}, that quotient has no support at $u=0$.
The cover $z^d=u$ is free of degree $d$ and is etale at nonzero $u$.
Lemma~\ref{lem:critical-ideal} therefore makes the contribution from
nonzero critical points of $R$ equal to
$d\,v(\mathcal D_{n,d})$.
Here the identifications are local; summing their finite lengths
is enough. The discriminant of a monic polynomial is, up to sign,
the determinant of multiplication by its derivative in the quotient
algebra. Its valuation is the length of that critical quotient.
The leading calculations above ensure the relevant modules are torsion.

Put $a_n=(\gamma)_{dn}$.
The root product formula for $R$ gives
\[
 \operatorname{Disc}_zR
       =\eta_{n,d}\,a_n^{d-1}D_{n,d}^{\,d},\qquad
           \eta_{n,d}\in\mathbb C^*,
\]
where the constant does not depend on $\gamma$.
At a parameter outside $\{0,-1,\ldots,-(dn-1)\}$,
$R(0)$ is a unit, so all critical support is nonzero and
$v(D_{n,d})=v(\mathcal D_{n,d})$.

At $\gamma_0=-m$, $0\le m<dn$, put $h=\lfloor m/d\rfloor$.
The specialized $R$ has at zero exactly the multiplicity
$L=d(n-h)$, since the last nonzero tilted coefficient is the one
with index $h$. Its constant coefficient has a simple zero in
$\gamma-\gamma_0$.
Split off the monic local factor of degree $L$ at zero; its remaining
factor is a unit there. Such a split follows by the coefficient
recursion for coprime monic factors modulo the maximal ideal.
All lower coefficients of the local factor are divisible by
$\varepsilon=\gamma-\gamma_0$, and its constant coefficient is
$\varepsilon$ times a unit.
After $\varepsilon=s^L$ and $z=sy$, division by $s^L$ leaves at
$s=0$ a polynomial $y^L+c$, $c\ne0$, with distinct roots.
Formal root lifting therefore makes every root difference have
$s$-valuation one. Its discriminant contribution has
$\varepsilon$-valuation $L-1$.
This counts the zero cluster without assuming that the nonzero
specialized roots are simple.
It follows that
\[
 (d-1)+d\,v(D_{n,d})=d\,v(\mathcal D_{n,d})+L-1,
 \qquad
 v(D_{n,d})-v(\mathcal D_{n,d})=n-1-\lfloor m/d\rfloor.
\]
Thus the rational function $D_{n,d}/\mathcal D_{n,d}$ has exactly
the finite divisor of
$\prod_{k=1}^{n-1}(\gamma)_{dk}$.
Over $\mathbb C$ their ratio is a nonzero constant.
The leading coefficients (L) make this constant
$(-1)^Nd^N\prod_{k=1}^nk^k$, proving Theorem~\ref{thm:discriminant}.
For $n=1$, both the fixed-degree discriminant and the empty
Hurwitz determinant are one, also directly verifying the formula.

\begin{corollary}\label{cor:top-nonnegative}
For every $n,d\ge1$ and every real $\gamma>0$,
$\mathcal D_{n,d}(\gamma)\ge0$, with strict inequality if and only
if $T_{n,\gamma,d}$ is squarefree.
At gap three and $n\ge2$ this gives, uniformly,
\[
 \Delta_{2n-2}(W)\ge0,\qquad
 \Delta_{2n-1}(W)\ge0,\qquad \Delta_{2n}(W)\ge0.
\]
All three are strictly positive together exactly when $T_{n,\gamma,3}$
is squarefree.
\end{corollary}
\begin{proof}
For $P_n(x)=\mathbb E(x+Y_\gamma^d)^n$, an even $n$ gives
$P_n(x)>0$ on the real axis. An odd $n$ gives
$P_n'(x)=n\mathbb E(x+Y_\gamma^d)^{n-1}>0$, so it has exactly
one real root and that root is simple.
If $P_n$ is squarefree, its other roots occur in
$\lfloor n/2\rfloor$ nonreal conjugate pairs, so the discriminant
has sign $(-1)^{\lfloor n/2\rfloor}=(-1)^N$.
If it has a repeated root, its discriminant is zero.
Hence $(-1)^ND_{n,d}\ge0$ in all cases.
Every other factor in (D) is positive, proving the first assertion.
At $d=3$, $B(0)>0$ and $W(0)>0$, so the two tail identities
in Theorem~\ref{thm:tail} give the remaining assertions.
\end{proof}

These are unbounded nonnegativity results, not a proof of strict
nonvanishing or of total nonnegativity of the entire generalized
Hurwitz matrix. At the physical shape $\gamma=n+\beta-1$,
strict positivity of the indicated determinant would close the
lower-one binomial family at that gap by the existing moment bridge.
General lower exponents require their separate boundary conditions.

\section{Negative minors at every gap}
\begin{theorem}\label{thm:negative}
For every integer $d\ge1$, at $b=1$, and every even integer $m\ge d$,
\[
 \det R(1,d)[\{1,\ldots,m\},\{0,\ldots,m-1\}]
                         =-\frac{m!}{d}<0.                 \tag{8}
\]
Consequently the physical auxiliary coefficient array is not totally
nonnegative at any gap $d\ge1$.
\end{theorem}
\begin{proof}
At $b=1$, (3) with $k=0$ has generating function
\[
 \sum_{n\ge0}(R^{-1})_{n,0}\frac{u^n}{n!}
 =\frac{P_d'(u)}{1-u}
 =\frac{(d+1)(1-u)^{d-1}-(1-u)^{-1}}d.
\]
The first term on the right is a polynomial of degree $d-1$.
Thus $(R^{-1})_{m,0}=-m!/d$ for every $m\ge d$.
The leading $(m+1)$-square block of $R$ is unit lower triangular.
The cofactor formula for its inverse gives
\[
 (R^{-1})_{m,0}=(-1)^m
       \det R[\{1,\ldots,m\},\{0,\ldots,m-1\}].
\]
For even $m$, this is (8). This is a uniform coefficient proof;
bounded replay examples are not used to extrapolate the statement.
\end{proof}

At $d=3,b=1,m=4$, the explicit minor is
\[
 \det\begin{pmatrix}
 3&1&0&0\\28&10&1&0\\440&160&21&1\\9720&3560&520&36
 \end{pmatrix}=-8.
\]
At arbitrary $b$, the same minor equals
\[
 (b-2)(b^3+4b^2-17b+20).
\]
The counterexample therefore occurs in the actual physical
arithmetic-progression family, not only in a generic inverse
polar-companion construction.

\section{The standard production matrix also fails}
Let $\Pi=R^{-1}\overline R$, where
$\overline R_{n,k}=R_{n+1,k}$. Then
\[
 A_d(u)=\frac1{P_d'(u)}=
       \frac d{(d+1)(1-u)^d-1},\qquad
 Z_{b,d}(u)=A_d'(u)+(b-2)\frac{A_d(u)}{1-u},
\]
and its exact entries are
\[
 \Pi_{i,j}=\frac{i!}{j!}
 \left([u^{i-j}]Z_{b,d}(u)+j[u^{i-j+1}]A_d(u)\right).       \tag{9}
\]
Coefficients at negative powers are zero. To derive (9), differentiate
the $j$th column generating function in (2), divide by $g$, and then
substitute $t=P_d(u)$. Equivalently the bivariate production generating
function is $e^{uv}(Z_{b,d}(u)+vA_d(u))$.

At the physical point $d=3,b=1$, its rows $2,3$ and columns $0,1$ give
\[
 \det\begin{pmatrix}166&62\\1874&698\end{pmatrix}=-320.
\]
These entries can also be reconstructed directly from the first five
rows of $R$ displayed above, including $R_{0,0}=1$.
Thus a proof that first establishes total nonnegativity of this
standard Riordan production matrix cannot apply here.

The failures in (8)--(9) concern $R$ and $\Pi$.
They do not refute total nonnegativity of the generalized Hurwitz matrix
of the gap-three Wronskian, the selected staircase determinants of that
Wronskian, or binomial Strong Factorial.
A valid positive-network proof for those objects must retain their
specific two-row Wronskian and cyclic-section structure instead of
assuming total nonnegativity of the whole auxiliary coefficient array.

The separate exact replay
\texttt{gap\_three\_hurwitz\_structure.py} checks symbolic shape
identities at stated finite gaps and orders, reconstructs the inverse
and production matrices, and verifies exact negative controls.
It also reconstructs the actual adjacent-row cyclic and sine-weight
identities symbolically and rejects an incorrect adjacent row.
Generic independent symbolic pairs and zero-constant or zero-minor
controls check the division-free tail identities in Theorem~\ref{thm:tail}.
Its receipt records those bounds and source hashes.
The all-gap identities, the cyclic Fitting presentation, the
multiplicity-sensitive discriminant identity and the infinite family
(8) have the proofs above. Additional controls check cyclic
presentations and reductions over fields and a nonreduced ring,
complex-shape adjoint identities, origin flags, zero-cluster widths,
leading constants and exact discriminant formulas at fixed indices. All computations are internal; independent replay,
formal verification, external review and priority claims are not supplied.
No all-index gap-three or unrestricted binomial Strong Factorial
closure is asserted.
\end{document}
